La \emph{devinette du menteur} est un jeu joué par deux joueurs $A$ et $B$. Les règles du jeu dépendent de
deux entiers strictement positifs $k$ et $n$, connus par chacun des deux joueurs.

Au début du jeu, le joueur $A$ choisit deux entiers $x$ et $N$ vérifiant $1 \leq x \leq N$. Le joueur $A$
garde $x$ secret et, honnêtement, communique $N$ au joueur $B$. Le joueur $B$ essaye d'obtenir des
informations concernant $x$ en posant au joueur $A$ des questions comme suit : pour chaque question, $B$
choisit un ensemble arbitraire d'entiers strictement positifs $S$ (éventuellement déjà choisi pour une
question antérieure) et demande à $A$ si $x$ appartient à $S$ ; le joueur $B$ peut poser autant de telles
questions qu'il le souhaite. Après chaque question, le joueur $A$ doit immédiatement répondre par \emph{oui}
ou \emph{non}, mais il a le droit de mentir autant de fois qu'il le souhaite ; la seule restriction étant que
parmi toutes $k + 1$ réponses consécutives, au moins l'une de ces réponses doit être la vérité.

Après que $B$ ait posé autant de questions qu'il le souhaite, il doit proposer un ensemble $X$ contenant au
plus $n$ entiers strictement positifs. Si $x$ appartient à $X$, alors $B$ gagne, sinon, il perd. Montrer que :
\begin{enumerate}
  \item Si $n \geq 2^k$, alors $B$ dispose d'une stratégie gagnante.
  \item Pour tout entier suffisamment grand $k$, il existe un entier $n \geq 1.99^k$ tel que $B$ ne dispose
  pas de stratégie gagnante.
\end{enumerate}
